At this point the target binary has been modified with:
\begin{itemize}
    \item A DT\_NEEDED entry has been added to \textit{.dynamic}, pointing to a shared object that contains 
        the forkserver declared as a constructor, alongside a RUNPATH that points to that shared object.
    \item A symbol representing the pointer to the \gls{SHM} has been inserted into \textit{.dynsym}.
    \item A new section that will be used as a personalized \gls{GOT} is added.
    \item A relocation that links the symbol to be resolved to the new section is added.
\end{itemize}
Now the \gls{IP}s must be identified so that they can subsequently be instrumented.
First, a distinction must be made between two types of instrumentation: 
\begin{itemize}
    \item An initial instrumentation is needed in order to initialize some variables needed for the \gls{SHM} update rule.
    \item All the other instrumentations perform the actual \gls{SHM} update.
\end{itemize}

In the first case, the \gls{IP} must be as close as possible to the \texttt{main} function of the binary.
Finding the main address is easy for both stripped and non-stripped binaries, so, once the main address is identified, 
the tool tries to find the first address after the function preamble.
This goal is not always achieved, because there is no unified grammar for the function prologue. 
It depends on how many registers the callee wants to save on the stack, how many variables the specific function needs \dots
In general, the first \gls{IP} lies in the first \gls{BB} of the main function.\\
All the other \gls{IP}s are found with a \gls{DFS} starting from the initial main \gls{BB}, searching among the 
successor \gls{BB}s of the current one.
Specifically, once a \gls{BB} is found and has not already been analyzed or discovered before, it is added to a queue.
\begin{algorithm}[ht]
    \caption{Instrumentation point search with DFS algorithm}
    \label{alg:ip_search_dfs}
    \begin{algorithmic}[1]
        \STATE $found \gets[\,\,]$ \COMMENT {\small \textcolor{green!60!black}{found is a list of discovered basic blocks. It contains both the basic blocks to be instrumented and the basic blocks for which instrumentation is not required}}
        \STATE $to\_instrument\gets[\,\,]$ \COMMENT{\small \textcolor{green!60!black}{this is the list of basic blocks to be instrumented. It is a subset of the ``found'' list}}
        \STATE $current=\texttt{main\_bb}$
        \STATE $stack\gets[\,\,]$
        \STATE $successors\gets \textsc{ExtractSuccessors}(current,found)$
        \FOR{$s \in successors$}
            \STATE $stack.\textsc{push}((current,s))$
        \ENDFOR
        \WHILE{$stack \text{ is not empty}$}
            \STATE $current\_edge\gets stack.\textsc{pop}()$
            \STATE $prev \gets current\_edge[\,0\,] , curr \gets current\_edge[\,1\,]$
            \IF{$prev.last\_instruction.\_opcode \neq ``call''$}
                \STATE $to\_instrument.\textsc{append}(curr)$  \COMMENT{\small \textcolor{green!60!black}{If the previous basic block ends with a call instruction, the current basic block is not instrumented}}
            \ENDIF
            \STATE $found.\textsc{append}(curr)$
            \STATE $successors\gets \textsc{ExtractSuccessors}(current,found)$
            \FOR{$s \in successors$}
                \STATE $stack.\textsc{push}((current,s))$ 
            \ENDFOR
        \ENDWHILE
    \end{algorithmic}
\end{algorithm}
The \gls{IP}s discovery procedure can be seen in \cref{alg:ip_search_dfs}.\\
Two different lists must be used in order to distinguish between found basic blocks and basic blocks to be instrumented.
For example, if a basic block ends with a call, it is ``useless'' to instrument the next one.
In fact, a \texttt{call} is like a branch that is always taken.
If an edge that is always taken is instrumented, this would not add any information for a fuzzer.
In fact, if, during fuzzing, the first basic block is discovered, the second will surely be discovered as well.
The situation is different for the \texttt{jmp} instruction. In fact, although this is also a branch that is always taken, it can be 
useful, for example in loops, to disambiguate between the backward edge (returning to the beginning of the loop) and the entry edge.
An example of this can be seen in \cref{fig:pow_example}.\\
The \textsc{ExtractSuccessors} function will find all the successors of a \gls{BB}, excluding :
\begin{itemize}
    \item The successors that are outside the \textit{``code''} region (so avoiding instrumenting \textit{.plt} stubs).
        For \textit{``code''} region is intended all the instruction of the main and other functions that are callable from the main.
        So avoid instrument functions in \textit{.fini\_array} and other not intereseting part of \textit{.text} section.
    \item The \gls{BB} that are already been found (and so are contained in ``found'' list)
    \item The \gls{BB} that are temporarily inside the stack data structure
\end{itemize} 
For the moment, it can be said that this function filters the successors of the current basic block, 
discarding those that it would make no sense to consider.
For example, if at a certain point the binary performs a call to a library function (e.g. \texttt{printf}), 
without excluding this case the program would also instrument the \textit{plt} stub of the library function.
\begin{figure}[ht]
    \centering
    \begin{subfigure}{0.8\textwidth}
        \centering
        \includegraphics[width=\textwidth,
            height=0.5\textheight
            ]{pow_main_example.png}
        \caption{Main CFG}
        \label{fig:pow_main_example}
    \end{subfigure}
    \begin{subfigure}{0.8\textwidth}
        \centering
        \includegraphics[width=\textwidth,
            height=0.35\textheight
            ]{pow_function_example.png}
        \caption{Pow function CFG}
        \label{fig:pow_function_example}
    \end{subfigure}
    \caption{Example of the CFG of a program that computes the power operation between two numbers}
    \label{fig:pow_example}
\end{figure}

In \cref{fig:pow_example} a disassembled \gls{CFG} obtained with Cutter \cite{cutter}, representing a program that computes the power operation, can be seen.
With respect to \cref{fig:pow_main_example}, the \gls{IP}s found are: 
\begin{itemize}
    \item \textbf{\textit{0x11c2}}: initial main \gls{BB}. This is where the initial instrumentation will be placed.
    \item \textbf{\textit{0x11f1}} (prev=\textit{0x11c2},next=\textit{0x11f1}) 
    \item \textbf{\textit{0x11db}} (prev=\textit{0x11c2},next=\textit{0x11db}) 
    \item \textbf{\textit{0x1253}} 
        \begin{itemize}
            \item (prev=\textit{0x11f1},next=\textit{0x1253})
            \item (prev=\textit{0x11db},next=\textit{0x1253}) 
        \end{itemize} 
    
\end{itemize}
Apart from the call to the \texttt{my\_pow} function, all the other functions are exported from libc, and so do not 
need to be instrumented.

With respect to \cref{fig:pow_function_example}, the \gls{IP}s found are:
\begin{itemize}
    \item \textbf{\textit{0x11b5}} 
        \begin{itemize}
            \item (prev=\textit{0x1189},next=\textit{0x11b5}) 
            \item (prev=\textit{0x11a7},next=\textit{0x11b5}) 
        \end{itemize}
    \item \textbf{\textit{0x11a7}} (prev=\textit{0x11b5},next=\textit{0x11a7}) 
    \item \textbf{\textit{0x11bd}} (prev=\textit{0x11b5},next=\textit{0x11bd}) 
\end{itemize}
The initial \gls{BB} of the function is not instrumented.
\clearpage